\section{Discussion and Conclusion}\label{sec:conclusion}
MINT-VLA improves held-out future-feature prediction, but the corrected
three-seed control experiment rejects action utility: A2-Abs and MINT-VLA are
$-3.94$ and $-8.58$ points below matched no-hint training, respectively.
Correct content also fails to beat fixed-checkpoint replacements.  Oracle stage
transitions, explicit next-subtask semantics, and privileged spatial targets do
not reverse the result.  A later visual state can identify progress without
specifying path, contact, arm assignment, or timing, so predictability and
representation alignment are not sufficient evidence of action sufficiency.

Preserving pretrained pathways does not yield a stable exception.  The
predictive adapter has a strong seed-1000 screen and a positive comparison with
one fixed A0, but independent matched A0s change the three seed effects to
+13.42, -5.50, and -2.08 points (mean +1.94; 95\% CI
[$-5.78$,$+12.75$]).  Its correct-versus-shuffled, route, mature-public, and
task-safety gates also fail.  Recurrence-aligned auxiliary training and
outcome-calibrated weighting provide further negative replication results.

The conclusion is deliberately scoped to the tested $\pi_{0.5}$ milestone and
predictive-adapter interfaces and the six-task RoboTwin panel.  It does not
rule out generative, action-conditioned, spatial, or planning-oriented world
models.  It does show that latent prediction metrics, one-seed gains, and
extra conditioning routes must not substitute for matched training seeds and
causal content controls.
